%%%
%%% Radical "⾃"
%%%
\section*{Radical 132: ``⾃''}\addcontentsline{toc}{section}{Radical 132: ⾃}\addcontentsline{loh}{figure}{\#\#\#\# 132: ⾃}

%%%%%%%%%% 自 %%%%%%%%%%
\subsection*{自}\addcontentsline{loh}{figure}{自}

\begin{Entry}{自}{6}{⾃}[Kangxi 132]
  \begin{Phonetics}{自}{zi4}[][HSK 4]
    \definition*{s.}{Sobrenome: Zi}
    \definition{adv.}{certamente; com certeza; é claro; naturalmente}
    \definition{prep.}{de; desde; a partir de; apresenta o ponto de partida, a fonte ou o horário de início do comportamento da ação, equivalente a 从 e 由}
    \definition{pron.}{si mesmo; próprio | próprio; indica que a ação é iniciada por e direcionada a si mesmo | por si mesmo; indica que a ação é autoiniciada e não é causada por uma força externa}
    \definition{v.}{iniciar}
  \seealsoref{从}{cong2}
  \seealsoref{由}{you2}
  \end{Phonetics}
\end{Entry}

\begin{Entry}{自力更生}{6,2,7,5}{⾃,⼒,⽈,⽣}
  \begin{Phonetics}{自力更生}{zi4li4-geng1sheng1}[][HSK 7-9]
    \definition{expr.}{autossuficiente; realize as coisas sem depender de forças externas, confiando na sua própria força}
  \synonymref{独立自主}{du2li4-zi4zhu3}
  \synonymref{艰苦奋斗}{jian1ku3-fen4dou4}
  \antonymref{不由自主}{bu4you2zi4zhu3}
  \end{Phonetics}
\end{Entry}

\begin{Entry}{自个儿}{6,3,2}{⾃,⼈,⼉}
  \begin{Phonetics}{自个儿}{zi4ge3r5}
    \definition{pron.}{Dialeto: a si mesmo, por si mesmo}
  \seealsoref{自各儿}{zi4ge3r5}
  \end{Phonetics}
\end{Entry}

\begin{Entry}{自卫}{6,3}{⾃,⼙}
  \begin{Phonetics}{自卫}{zi4wei4}[][HSK 7-9]
    \definition{s.}{autodefesa}
    \definition{v.}{defender"-se}
  \antonymref{伤害}{shang1hai4}
  \end{Phonetics}
\end{Entry}

\begin{Entry}{自己}{6,3}{⾃,⼰}
  \begin{Phonetics}{自己}{zi4ji3}[][HSK 2]
    \definition{pron.}{a si próprio; a si mesmo; refere"-se ao substantivo ou pronome precedente (enfatiza principalmente que não é devido a forças externas)}
  \end{Phonetics}
\end{Entry}

\begin{Entry}{自己动手}{6,3,6,4}{⾃,⼰,⼒,⼿}
  \begin{Phonetics}{自己动手}{zi4ji3 dong4shou3}
    \definition{v.}{usar as próprias mãos; fazer por si mesmo | fazer algo por si mesmo | servir"-se de}
  \end{Phonetics}
\end{Entry}

\begin{Entry}{自从}{6,4}{⾃,⼈}
  \begin{Phonetics}{自从}{zi4cong2}[][HSK 3]
    \definition{prep.}{de; desde; a partir de; referir"-se a um momento ou evento específico no passado}
  \end{Phonetics}
\end{Entry}

\begin{Entry}{自以为是}{6,4,4,9}{⾃,⼈,⼂,⽇}
  \begin{Phonetics}{自以为是}{zi4yi3wei2shi4}[][HSK 7-9]
    \definition{expr.}{considerar"-se (sempre) correto; considerar"-se infalível; ter opinião própria; não admitir erros; ele acredita que seus pontos de vista e ações estão corretos e não aceita as opiniões dos outros}
  \antonymref{不耻下问}{bu4chi3-xia4wen4}
  \end{Phonetics}
\end{Entry}

\begin{Entry}{自主}{6,5}{⾃,⼂}
  \begin{Phonetics}{自主}{zi4zhu3}[][HSK 3]
    \definition{v.}{agir por conta própria; decidir por si mesmo; manter a iniciativa em suas próprias mãos; tomar suas próprias decisões}
  \end{Phonetics}
\end{Entry}

\begin{Entry}{自发}{6,5}{⾃,⼜}
  \begin{Phonetics}{自发}{zi4fa1}[][HSK 7-9]
    \definition{adj.}{espontâneo; que ocorre naturalmente sem influência externa}
  \synonymref{自觉}{zi4jue2}
  \synonymref{自愿}{zi4yuan4}
  \antonymref{强制}{qiang2zhi4}
  \end{Phonetics}
\end{Entry}

\begin{Entry}{自号}{6,5}{⾃,⼝}
  \begin{Phonetics}{自号}{zi4hao4}
    \definition{s.}{autoproclamado; auto"-número}
  \end{Phonetics}
\end{Entry}

\begin{Entry}{自由}{6,5}{⾃,⽥}
  \begin{Phonetics}{自由}{zi4you2}[][HSK 2]
    \definition{adj.}{livre; irrestrito}
    \definition[个]{s.}{liberdade; o direito de agir de acordo com a própria vontade dentro do âmbito da lei | liberdade; filosoficamente, liberdade é definida como o processo de as pessoas reconhecerem as leis que governam o desenvolvimento das coisas e aplicá"-las conscientemente na prática}
  \end{Phonetics}
\end{Entry}

\begin{Entry}{自由自在}{6,5,6,6}{⾃,⽥,⾃,⼟}
  \begin{Phonetics}{自由自在}{zi4you2-zi4zai4}[][HSK 7-9]
    \definition{expr.}{livremente; à vontade; sem restrições, de espírito livre, pacífico e descontraído}
  \antonymref{身不由己}{shen1bu4you2ji3}
  \end{Phonetics}
\end{Entry}

\begin{Entry}{自由泳}{6,5,8}{⾃,⽥,⽔}
  \begin{Phonetics}{自由泳}{zi4you2yong3}
    \definition{s.}{estilo livre (natação); \emph{crawl}; \emph{prone"-crawl}}
  \end{Phonetics}
\end{Entry}

\begin{Entry}{自立}{6,5}{⾃,⽴}
  \begin{Phonetics}{自立}{zi4li4}[][HSK 7-9]
    \definition{adj.}{independente; autossuficiente; autossustentável}
    \definition{v.}{ficar de pé por conta própria;  sustentar"-se; viver do próprio trabalho, sem depender de outros}
  \synonymref{独立}{du2li4}
  \synonymref{自主}{zi4zhu3}
  \synonymref{自助}{zi4zhu4}
  \antonymref{寄托}{ji4tuo1}
  \antonymref{依靠}{yi1kao4}
  \antonymref{依赖}{yi1lai4}
  \end{Phonetics}
\end{Entry}

\begin{Entry}{自动}{6,6}{⾃,⼒}
  \begin{Phonetics}{自动}{zi4dong4}[][HSK 3]
    \definition{adj.}{automático; auto"-atuante; uso de dispositivos mecânicos, elétricos, etc, para funcionar automaticamente, sem necessidade de controle humano}
    \definition{adv.}{voluntariamente; por vontade própria; por iniciativa própria | automaticamente; espontaneamente; refere"-se a movimentos, mudanças, etc., que não são causados pela ação humana, mas sim pelo próprio objeto}
  \end{Phonetics}
\end{Entry}

\begin{Entry}{自动化}{6,6,4}{⾃,⼒,⼔}
  \begin{Phonetics}{自动化}{zi4dong4hua4}
    \definition{s.}{automatizar}
    \definition{s.}{robotização; automação; automatização}
  \synonymref{机器人}{ji1qi4ren2}
  \end{Phonetics}
\end{Entry}

\begin{Entry}{自各儿}{6,6,2}{⾃,⼝,⼉}
  \begin{Phonetics}{自各儿}{zi4ge3r5}
    \definition{pron.}{Dialeto: si mesmo; por si mesmo}
  \end{Phonetics}
\end{Entry}

\begin{Entry}{自在}{6,6}{⾃,⼟}
  \begin{Phonetics}{自在}{zi4zai5}[][HSK 6]
    \definition{adj.}{livre; irrestrito; liberto}
  \end{Phonetics}
\end{Entry}

\begin{Entry}{自如}{6,6}{⾃,⼥}
  \begin{Phonetics}{自如}{zi4ru2}[][HSK 7-9]
    \definition{adj.}{suave; livre; as atividades ou operações não são prejudicadas | normal; calmo e natural}
  \synonymref{自在}{zi4zai5}
  \end{Phonetics}
\end{Entry}

\begin{Entry}{自杀}{6,6}{⾃,⽊}
  \begin{Phonetics}{自杀}{zi4sha1}[][HSK 5]
    \definition{s.}{suicídio; auto"-assassinato; auto"-sacrifício}
    \definition{v.}{cometer suicídio; tentar suicídio; matar"-se}
  \end{Phonetics}
\end{Entry}

\begin{Entry}{自行}{6,6}{⾃,⾏}
  \begin{Phonetics}{自行}{zi4xing2}[][HSK 7-9]
    \definition{adv.}{autonomamente; por si só; faça você mesmo (sem depender de outros ou por meio deles) | por si mesmo; por iniciativa própria; voluntariamente; ocorre ou surge naturalmente, sem intervenção humana}
  \synonymref{独立}{du2li4}
  \synonymref{自动}{zi4dong4}
  \synonymref{自己}{zi4ji3}
  \synonymref{自主}{zi4zhu3}
  \end{Phonetics}
\end{Entry}

\begin{Entry}{自行车}{6,6,4}{⾃,⾏,⾞}
  \begin{Phonetics}{自行车}{zi4xing2che1}[][HSK 2]
    \definition[辆]{s.}{bicicleta; um veículo de duas rodas que é impulsionado para a frente com os pedais}
  \end{Phonetics}
\end{Entry}

\begin{Entry}{自行车架}{6,6,4,9}{⾃,⾏,⾞,⽊}
  \begin{Phonetics}{自行车架}{zi4xing2che1 jia4}
    \definition{s.}{quadro de bicicleta | suporte (ou \emph{rack}) para bicicletas}
  \end{Phonetics}
\end{Entry}

\begin{Entry}{自行车馆}{6,6,4,11}{⾃,⾏,⾞,⾷}
  \begin{Phonetics}{自行车馆}{zi4xing2che1 guan3}
    \definition{s.}{estádio de ciclismo | velódromo}
  \end{Phonetics}
\end{Entry}

\begin{Entry}{自行车赛}{6,6,4,14}{⾃,⾏,⾞,⾙}
  \begin{Phonetics}{自行车赛}{zi4xing2che1sai4}
    \definition{s.}{corrida de bicicleta}
  \end{Phonetics}
\end{Entry}

\begin{Entry}{自负}{6,6}{⾃,⾙}
  \begin{Phonetics}{自负}{zi4fu4}[][HSK 7-9]
    \definition{adj.}{presunçoso; ter alta opinião de si mesmo; ter uma autoestima excessiva}
    \definition{v.}{ser responsável pelos próprios atos; assumir a responsabilidade você mesmo}
  \synonymref{自信}{zi4xin4}
  \antonymref{谦虚}{qian1xu1}
  \antonymref{自卑}{zi4bei1}
  \end{Phonetics}
\end{Entry}

\begin{Entry}{自助}{6,7}{⾃,⼒}
  \begin{Phonetics}{自助}{zi4zhu4}[][HSK 7-9]
    \definition{s.}{autosserviço}
    \definition{v.}{depender de si mesmo; recorrer à autoajuda; apoiar/ajudar/servir a si mesmo}
  \synonymref{自立}{zi4li4}
  \synonymref{自主}{zi4zhu3}
  \end{Phonetics}
\end{Entry}

\begin{Entry}{自我}{6,7}{⾃,⼽}
  \begin{Phonetics}{自我}{zi4wo3}[][HSK 6]
    \definition{pref.}{auto-}
    \definition{pron.}{a si mesmo; eu próprio; geralmente usado antes de verbos dissílabos para indicar que a ação é realizada por alguém e dirigida a si mesmo | indicar o próprio caráter diferente dos outros; refere"-se às próprias características, personalidade, hobbies, etc.}
  \end{Phonetics}
\end{Entry}

\begin{Entry}{自我介绍}{6,7,4,8}{⾃,⼽,⼈,⽷}
  \begin{Phonetics}{自我介绍}{zi4wo3jie4shao4}
    \definition{s.}{autoapresentação; apresentar"-se; apresente seu histórico familiar e sua experiência educacional.}
  \end{Phonetics}
\end{Entry}

\begin{Entry}{自我安慰}{6,7,6,15}{⾃,⼽,⼧,⼼}
  \begin{Phonetics}{自我安慰}{zi4wo3'an1wei4}
    \definition{s.}{autoconsolação}
    \definition{v.}{confortar a si mesmo | consolar a si mesmo | tranquilizar a si mesmo}
  \end{Phonetics}
\end{Entry}

\begin{Entry}{自我防卫}{6,7,6,3}{⾃,⼽,⾩,⼙}
  \begin{Phonetics}{自我防卫}{zi4wo3 fang2wei4}
    \definition{s.}{defesa pessoal | auto"-defesa}
  \end{Phonetics}
\end{Entry}

\begin{Entry}{自我吹嘘}{6,7,7,14}{⾃,⼽,⼝,⼝}
  \begin{Phonetics}{自我吹嘘}{zi4wo3 chui1xu1}
    \definition{expr.}{gabar"-se}
    \definition{s.}{auto"-ostentação; autoglorificação}
  \end{Phonetics}
\end{Entry}

\begin{Entry}{自我批评}{6,7,7,7}{⾃,⼽,⼿,⾔}
  \begin{Phonetics}{自我批评}{zi4wo3pi1ping2}
    \definition{s.}{autocrítica}
  \end{Phonetics}
\end{Entry}

\begin{Entry}{自我实现}{6,7,8,8}{⾃,⼽,⼧,⾒}
  \begin{Phonetics}{自我实现}{zi4wo3shi2xian4}
    \definition{s.}{Psicologia: auto"-realização}
  \end{Phonetics}
\end{Entry}

\begin{Entry}{自我的人}{6,7,8,2}{⾃,⼽,⽩,⼈}
  \begin{Phonetics}{自我的人}{zi4wo3 de5 ren2}
    \definition{s.}{(minha, sua) própria pessoa | (afirmar) a própria personalidade}
  \end{Phonetics}
\end{Entry}

\begin{Entry}{自我保存}{6,7,9,6}{⾃,⼽,⼈,⼦}
  \begin{Phonetics}{自我保存}{zi4wo3 bao3cun2}
    \definition{v.}{autopreservação}
  \end{Phonetics}
\end{Entry}

\begin{Entry}{自我陶醉}{6,7,10,15}{⾃,⼽,⾩,⾣}
  \begin{Phonetics}{自我陶醉}{zi4wo3tao2zui4}
    \definition{s.}{extrema complacência; embriagado de si mesmo; muito satisfeito consigo mesmo; narcisista; autoimpregnado; autossatisfeito}
  \synonymref{如醉如痴}{ru2zui4-ru2chi1}
  \end{Phonetics}
\end{Entry}

\begin{Entry}{自我催眠}{6,7,13,10}{⾃,⼽,⼈,⽬}
  \begin{Phonetics}{自我催眠}{zi4wo3 cui1mian2}
    \definition{s.}{autohipnotismo; autohipnose; idiohipnotismo}
  \end{Phonetics}
\end{Entry}

\begin{Entry}{自我意识}{6,7,13,7}{⾃,⼽,⼼,⾔}
  \begin{Phonetics}{自我意识}{zi4wo3yi4shi2}
    \definition{s.}{autoconsciência; auto"-consciente}
  \end{Phonetics}
\end{Entry}

\begin{Entry}{自我解嘲}{6,7,13,15}{⾃,⼽,⾓,⼝}
  \begin{Phonetics}{自我解嘲}{zi4wo3jie3chao2}
    \definition{s.}{autodepreciação; referir"-se às próprias fraquezas ou falhas com humor autodepreciativo}
    \definition{v.}{encontrar desculpas para; consolar"-se}
  \end{Phonetics}
\end{Entry}

\begin{Entry}{自来水}{6,7,4}{⾃,⽊,⽔}
  \begin{Phonetics}{自来水}{zi4lai2shui3}[][HSK 6]
    \definition{s.}{água da torneira; água corrente; água purificada e desinfetada fornecida por sistema hidráulico através de tubulações | equipamentos para transporte de água natural tratada}
  \end{Phonetics}
\end{Entry}

\begin{Entry}{自私}{6,7}{⾃,⽲}
  \begin{Phonetics}{自私}{zi4si1}[][HSK 7-9]
    \definition{adj.}{egoísta; egocêntrico; interesseiro; que só se importam com os próprios interesses e ignoram os outros}
  \antonymref{大方}{da4fang5}
  \antonymref{慷慨}{kang1kai3}
  \antonymref{无私}{wu2si1}
  \end{Phonetics}
\end{Entry}

\begin{Entry}{自私自利}{6,7,6,7}{⾃,⽲,⾃,⼑}
  \begin{Phonetics}{自私自利}{zi4si1-zi4li4}[][HSK 7-9]
    \definition[个]{expr.}{egoísta; pensar apenas em si mesmo; preocupar"-se apenas com os próprios interesses}
  \synonymref{损人利己}{sun3ren2-li4ji3}
  \antonymref{大公无私}{da4gong1-wu2si1}
  \end{Phonetics}
\end{Entry}

\begin{Entry}{自言自语}{6,7,6,9}{⾃,⾔,⾃,⾔}
  \begin{Phonetics}{自言自语}{zi4yan2-zi4yu3}[][HSK 6]
    \definition{expr.}{falando sozinho; falar consigo mesmo; pensar em voz alta; solilóquio}
  \end{Phonetics}
\end{Entry}

\begin{Entry}{自身}{6,7}{⾃,⾝}
  \begin{Phonetics}{自身}{zi4shen1}[][HSK 3]
    \definition{pron.}{eu mesmo (enfatizando que não é outra pessoa ou outra coisa)}
  \end{Phonetics}
\end{Entry}

\begin{Entry}{自卑}{6,8}{⾃,⼗}
  \begin{Phonetics}{自卑}{zi4bei1}[][HSK 7-9]
    \definition{adj.}{abjeto; que se sente inferior; auto"-humilhação}
  \synonymref{惭愧}{can2kui4}
  \synonymref{散心}{san4/xin1}
  \antonymref{谦逊}{qian1xun4}
  \antonymref{虚心}{xu1xin1}
  \antonymref{自负}{zi4fu4}
  \antonymref{自豪}{zi4hao2}
  \antonymref{自信}{zi4xin4}
  \antonymref{自尊}{zi4zun1}
  \end{Phonetics}
\end{Entry}

\begin{Entry}{自始至终}{6,8,6,8}{⾃,⼥,⾄,⽷}
  \begin{Phonetics}{自始至终}{zi4shi3-zhi4zhong1}[][HSK 7-9]
    \definition{expr.}{``Do começo ao fim.''; do início ao fim}
  \synonymref{一如既往}{yi4ru2-ji4wang3}
  \antonymref{半途而废}{ban4tu2'er2fei4}
  \antonymref{朝三暮四}{zhao1san1-mu4si4}
  \end{Phonetics}
\end{Entry}

\begin{Entry}{自学}{6,8}{⾃,⼦}
  \begin{Phonetics}{自学}{zi4xue2}[][HSK 6]
    \definition{s.}{auto"-estudo; autodidata; autoaprendizagem}
    \definition{v.}{estudar por conta própria; estudar de forma independente; ensinar a si mesmo}
  \end{Phonetics}
\end{Entry}

\begin{Entry}{自责}{6,8}{⾃,⾙}
  \begin{Phonetics}{自责}{zi4ze2}[][HSK 7-9]
    \definition{v.}{culpar a si mesmo; repreender a si mesmo}
  \end{Phonetics}
\end{Entry}

\begin{Entry}{自信}{6,9}{⾃,⼈}
  \begin{Phonetics}{自信}{zi4xin4}[][HSK 4]
    \definition{adj.}{confiante; descreve a crença em suas próprias habilidades, decisões, etc., tendo confiança em si mesmo}
    \definition[份,种]{s.}{autoconfiança; confiança em si mesmo}
    \definition{v.}{acreditar em si mesmo}
  \end{Phonetics}
\end{Entry}

\begin{Entry}{自信心}{6,9,4}{⾃,⼈,⼼}
  \begin{Phonetics}{自信心}{zi4xin4xin1}[][HSK 7-9]
    \definition{s.}{autoconfiança; um estado de espírito que acredita ter a capacidade de realizar certos desejos}
  \synonymref{自尊心}{zi4zun1xin1}
  \end{Phonetics}
\end{Entry}

\begin{Entry}{自相矛盾}{6,9,5,9}{⾃,⽬,⽭,⽬}
  \begin{Phonetics}{自相矛盾}{zi4xiang1-mao2dun4}[][HSK 7-9]
    \definition{expr.}{inconsistente; autocontradição; antinomias; uma situação em que as palavras, ações ou opiniões de uma pessoa são inconsistentes, levando a uma autocontradição; contradizer"-se; autocontraditório}
  \synonymref{格格不入}{ge2ge2-bu2ru4}
  \end{Phonetics}
\end{Entry}

\begin{Entry}{自觉}{6,9}{⾃,⾒}
  \begin{Phonetics}{自觉}{zi4jue2}[][HSK 3]
    \definition{adj.}{autoconsciente; de livre e espontânea vontade; controlar o próprio comportamento e agir por iniciativa própria}
    \definition{v.}{estar ciente de}
  \end{Phonetics}
\end{Entry}

\begin{Entry}{自费}{6,9}{⾃,⾙}
  \begin{Phonetics}{自费}{zi4fei4}[][HSK 7-9]
    \definition{v.}{pagar por conta própria; ser às custas do próprio bolso; ser autofinanciado}
  \antonymref{公费}{gong1fei4}
  \end{Phonetics}
\end{Entry}

\begin{Entry}{自称}{6,10}{⾃,⽲}
  \begin{Phonetics}{自称}{zi4cheng1}[][HSK 7-9]
    \definition[个]{s.}{autodenominar"-se; afirmar ser; professar; autoproclamar"-se | declarar"-se como sendo; alegar | chamar"-se (estilizar"-se)}
  \seealsoref{自号}{zi4hao4}
  \end{Phonetics}
\end{Entry}

\begin{Entry}{自救}{6,11}{⾃,⽁}
  \begin{Phonetics}{自救}{zi4jiu4}
    \definition{s.}{autoajuda}
    \definition{v.}{salvar"-se; prover"-se e ajudar"-se}
  \end{Phonetics}
\end{Entry}

\begin{Entry}{自理}{6,11}{⾃,⽟}
  \begin{Phonetics}{自理}{zi4li3}[][HSK 7-9]
    \definition{v.}{prover o próprio sustento; pagar você mesmo | cuidar de si mesmo; cuidar dos seus próprios assuntos; ter autocuidado}
  \end{Phonetics}
\end{Entry}

\begin{Entry}{自尊}{6,12}{⾃,⼨}
  \begin{Phonetics}{自尊}{zi4zun1}[][HSK 7-9]
    \definition{s.}{respeito próprio; autoestima; orgulho adequado; respeite"-se; não se humilhe perante os outros nem tolere discriminação ou insultos}
  \synonymref{自负}{zi4fu4}
  \synonymref{自豪}{zi4hao2}
  \synonymref{自信}{zi4xin4}
  \antonymref{自卑}{zi4bei1}
  \end{Phonetics}
\end{Entry}

\begin{Entry}{自尊心}{6,12,4}{⾃,⼨,⼼}
  \begin{Phonetics}{自尊心}{zi4zun1xin1}[][HSK 7-9]
    \definition{s.}{autoestima; senso de autoestima; a ideia de respeitar a si mesmo e não ser submisso aos outros}
  \synonymref{自信心}{zi4xin4xin1}
  \end{Phonetics}
\end{Entry}

\begin{Entry}{自强不息}{6,12,4,10}{⾃,⼸,⼀,⼼}
  \begin{Phonetics}{自强不息}{zi4qiang2-bu4xi1}[][HSK 7-9]
    \definition{expr.}{busque o aprimoramento pessoal; esforçar"-se constantemente para se tornar mais forte; fazer esforços incessantes para se aprimorar; autoaperfeiçoamento; lutar incessantemente}
  \synonymref{发愤图强}{fa1fen4-tu2qiang2}
  \synonymref{发奋图强}{fa1fen4-tu2qiang2}
  \end{Phonetics}
\end{Entry}

\begin{Entry}{自然}{6,12}{⾃,⽕}
  \begin{Phonetics}{自然}{zi4ran5}[][HSK 3]
    \definition{adj.}{natural; no curso normal dos eventos; formado ou desenvolvido sem intervenção humana; algo que se desenvolve livremente}
    \definition{adv.}{naturalmente; definitivamente; certamente, isso significa que, de acordo com a lógica, deve ser assim}
    \definition{conj.}{usado para ligar duas frases, com a segunda introduzindo informações adicionais ou adversativas; indica explicação complementar ou uma mudança de significado}
    \definition{s.}{natureza; mundo natural; tudo o que não foi criado pelo ser humano}
  \end{Phonetics}
\end{Entry}

\begin{Entry}{自然而然}{6,12,6,12}{⾃,⽕,⽽,⽕}
  \begin{Phonetics}{自然而然}{zi4ran2'er2ran2}[][HSK 7-9]
    \definition{expr.}{naturalmente; automaticamente; de si mesmo; isso acontece sem a ação de uma força externa}
  \synonymref{顺其自然}{shun4qi2zi4ran2}
  \end{Phonetics}
\end{Entry}

\begin{Entry}{自然界}{6,12,9}{⾃,⽕,⽥}
  \begin{Phonetics}{自然界}{zi4ran2jie4}[][HSK 7-9]
    \definition{s.}{mundo natural; natureza}
  \synonymref{大自然}{da4zi4ran2}
  \end{Phonetics}
\end{Entry}

\begin{Entry}{自愿}{6,14}{⾃,⽕}
  \begin{Phonetics}{自愿}{zi4yuan4}[][HSK 5]
    \definition{adv.}{voluntariamente; por iniciativa própria; por vontade própria}
    \definition{s.}{voluntário}
  \end{Phonetics}
\end{Entry}

\begin{Entry}{自豪}{6,14}{⾃,⾗}
  \begin{Phonetics}{自豪}{zi4hao2}[][HSK 5]
    \definition{adj.}{orgulhar"-se de; ter orgulho de; sentir"-se honrado por possuir qualidades excelentes ou ter alcançado grandes conquistas, seja por si mesmo ou por um grupo ou indivíduo relacionado a si}
  \end{Phonetics}
\end{Entry}

\begin{Entry}{自燃}{6,16}{⾃,⽕}
  \begin{Phonetics}{自燃}{zi4ran2}
    \definition{s.}{combustão espontânea; autoignição; puroforicidade; substâncias podem entrar em combustão espontânea por meio de oxidação lenta no ar; por exemplo, o fósforo branco pode inflamar"-se espontaneamente, e grandes pilhas de carvão, algodão e feno também podem entrar em combustão espontânea em condições de pouca ventilação}
  \end{Phonetics}
\end{Entry}

%%%%%%%%%% 臭 %%%%%%%%%%
\subsection*{臭}\addcontentsline{loh}{figure}{臭}

\begin{Entry}{臭}{10}{⾃}
  \begin{Phonetics}{臭}{chou4}[][HSK 5]
    \definition{adj.}{sujo; malcheiroso; fedorento; contrário de 香 | repugnante; nojento; repulsivo | ruim; pobre; péssimo}
    \definition{adv.}{severamente; firmemente}
    \definition{v.}{falhar em detonar (bala)}
  \seealsoref{香}{xiang1}
  \end{Phonetics}
  \begin{Phonetics}{臭}{xiu4}
    \definition{s.}{odor; cheiro}
    \definition{v.}{cheirar; farejar; o mesmo que 嗅}
  \seealsoref{嗅}{xiu4}
  \end{Phonetics}
\end{Entry}

\begin{Entry}{臭气}{10,4}{⾃,⽓}
  \begin{Phonetics}{臭气}{chou4qi4}
    \definition{s.}{mau cheiro (ou odor ofensivo); fedor; gases que emitem um odor desagradável; por exemplo, a amônia | mau cheiro (ofensa); sabor desagradável; fedor; umidade fétida | mau cheiro; odor fétido}
  \antonymref{香气}{xiang1qi4}
  \end{Phonetics}
\end{Entry}

%%%%% EOF %%%%%

